\begin{solapartat}
\figura[0.45]{sol_fig1.png}

\punts{0,6} La longitud d'ona de l'harmònic fonamental és el doble de la longitud de la cavitat. Dins de la cavitat tenim mitja longitud d'ona (vegeu figura).
\[ \lambda = 2L = 64\ \mathrm{cm} = 0,64\ \mathrm{m} \]
L'harmònic fonamental té dos nodes (extrems) i un ventre (al centre).

\punts{0,65} En el cas del segon harmònic tenim tres nodes i dos ventres:

\figura[0.45]{sol_fig2.png}
\end{solapartat}

\begin{solapartat}
\punts{0,65} La freqüència no canvia encara que canviem de medi de transmissió $f = 196\ \mathrm{Hz}$.

\punts{0,6} $\lambda = \dfrac{v}{f} = \dfrac{340}{196} = 1,73\ \mathrm{m}$
\end{solapartat}
